Currently, due to time constraints, some RSSI values are missing from the dataset, and some must be repeated, as evident from figures \ref{fig:advertisements-distance} and \ref{fig:rssi-density}: the 50 and 100 meter plots are missing, the 80, 90, and 110 meter plots are jittery - some of which due to having an abnornally low received packet count, which indicates the experiment should be repeated in future work.
The RSSI values we have provided for the ranges that are complete and stable, are mostly sufficient to derive the necessary path loss parameters for the simulation, as they exhibit the expected power-law distribution and are mostly consistent.
Figure \ref{fig:latency-percentiles} interestingly shows that latency increases with distance, which confirms the path-loss model, but exhibits an unexpected result: we expected the 120m radius simulation result to have the highest latency on average, as it created the most messages and suffers the most from the epidemic routing's flooding of messages, but it was dominated by 50-70m ranges, perhaps due to the missing RSSI values at those distances.
The centralization metrics plot (Figure \ref{fig:centralization-metrics}) summarizes the simulation results, and proves that L0 loss is indeed not enough to characterise centralization score: it has 1.0 for most values.
The Gini index and the Centralization score S plots show a positive relation: as the Gini index increases, the S score increases as well, but the curves that dominate the graph are the not the same: the Gini index is dominated by the 40m range, while the S score is dominated by the 120m range.